%% \BibTeX command to typeset BibTeX logo in the docs
\AtBeginDocument{%
    \providecommand\BibTeX{{%
        \normalfont B\kern-0.5em{\scshape i\kern-0.25em b}\kern-0.8em\TeX}}}


%%%%%%%%%%%%%%%%%%%%%%%%%%%%%%%%%%%%%
%             CUSTOM                %
%%%%%%%%%%%%%%%%%%%%%%%%%%%%%%%%%%%%%
\usepackage[utf8]{inputenc}
\usepackage[T1]{fontenc}

\usepackage{soul}
\usepackage{subcaption}
\usepackage{lipsum}
\usepackage{multicol}% http://ctan.org/pkg/multicols


\newcommand{\fig}[1]{Fig.~\ref{#1}}
\newcommand{\listing}[1]{Listing~\ref{#1}}
\newcommand{\alg}[1]{Algorithm~\ref{#1}}
\newcommand{\tbl}[1]{Table~\ref{#1}}
\newcommand{\annex}[1]{Annex~\ref{#1}}
\newcommand{\sect}[1]{Section~\ref{#1}}

\usepackage{multirow}

%%%%% AVOID PAGE BREAKS %%%%%
\newenvironment{absolutelynopagebreak}
  {\par\nobreak\vfil\penalty0\vfilneg
   \vtop\bgroup}
  {\par\xdef\tpd{\the\prevdepth}\egroup
   \prevdepth=\tpd}

%%%%%% COMMENTAIRES %%%%%%
\usepackage[]{ifthen}
\newcommand\showComments{1} % 0 = False, 1 = True
\newcommand{\romain}[1]{\ifthenelse{\equal{\showComments}{0}}{}{\textcolor{blue}{Romain: #1}}}
\newcommand{\clement}[1]{\ifthenelse{\equal{\showComments}{0}}{}{\textcolor{purple}{Clément: #1}}}
\newcommand{\adrien}[1]{\ifthenelse{\equal{\showComments}{0}}{}{\textcolor{red}{Adrien: #1}}}
\newcommand{\patrick}[1]{\ifthenelse{\equal{\showComments}{0}}{}{\textcolor{gray}{Patrick: #1}}}
\newcommand{\thierry}[1]{\ifthenelse{\equal{\showComments}{0}}{}{\textcolor{orange}{Thierry: #1}}}

%% \code %%
\makeatletter
\DeclareRobustCommand*{\escapeus}[1]{%
    \begingroup\@activeus\scantokens{#1\endinput}\endgroup}
    \begingroup\lccode`\~=`\_\relax
    \lowercase{\endgroup\def\@activeus{\catcode`\_=\active \let~\_}}
    \makeatother
\def\code#1{\texttt{\escapeus{#1}}}

%% \ceil and \floor %%
\usepackage{mathtools}
\DeclarePairedDelimiter\ceil{\lceil}{\rceil}
\DeclarePairedDelimiter\floor{\lfloor}{\rfloor}

% align figure left/right
\usepackage[export]{adjustbox}

%%%%%%%%%%%%%%%%%%%%%%%%%%%%%%%%%%%%%%%%%%%%%%%%%%%%%%%%%%%%%
%               LISTINGS                                    %
%%%%%%%%%%%%%%%%%%%%%%%%%%%%%%%%%%%%%%%%%%%%%%%%%%%%%%%%%%%%%
\usepackage{listings}
\usepackage{color}
\definecolor{bluekeywords}{rgb}{0.13,0.13,1}
\definecolor{greencomments}{rgb}{0,0.5,0}
\definecolor{redstrings}{rgb}{0.9,0,0}

\usepackage{listings}
% \lstset{language=C,
%     frame=single,
%     showspaces=false,
%     showtabs=false,
%     breaklines=true,
%     numbers=left,
%     showstringspaces=false,
%     breakatwhitespace=true,
%     escapeinside={(*@}{@*)},
%     commentstyle=\color{greencomments},
%     keywordstyle=\color{bluekeywords},
%     stringstyle=\color{redstrings},
%     basicstyle=\ttfamily
% }

\definecolor{mCyan}{rgb}{0.17,0.69,0.79}
\definecolor{mGreen}{rgb}{0.08,0.70,0.15}
\definecolor{mBrown}{rgb}{0.6,0.6,0.06}
\definecolor{mBlue}{rgb}{0.39,0.39,1.0}

\lstdefinestyle{CStyle}{
    framextopmargin=50pt,
    frame=bottomline,
    basicstyle=\footnotesize\ttfamily,
    breaklines=true,
    columns=flexible,
    stepnumber=1,
    numbers=left,
    tabsize=4,
    language=C,
    keywordstyle=\bfseries\color{mGreen},
    commentstyle=\color{mCyan},
    numberstyle=\color{mBrown},
    escapeinside={<@}{@>},
    morekeywords={_Thread_local}
}


%%%%%%%%%%%%%%%%%%%%%%%%%%%%%%%%%%%%%%%%%%%%%%%%%%%%%%%%%%%%%%%%%%%%%%%%%%%%%%%%%%%
%               PSEUDO-CODE REQUIRED PACKAGES                                     %
%%%%%%%%%%%%%%%%%%%%%%%%%%%%%%%%%%%%%%%%%%%%%%%%%%%%%%%%%%%%%%%%%%%%%%%%%%%%%%%%%%%

\usepackage{amsmath}
\usepackage{algorithm}
\usepackage[normalem]{ulem}
%\usepackage[ruled,vlined,linesnumbered]{algorithm2e}
\usepackage[noend]{algpseudocode}
\usepackage{etoolbox}
\makeatletter

\algnewcommand\algorithmicforeach{\textbf{for each}}
\algdef{S}[FOR]{ForEach}[1]{\algorithmicforeach\ #1\ \algorithmicdo}

% Defines custom \On ... \EndOn block
\algblockdefx[LOCK]{Lock}{Unlock}[1]{\textbf{lock}~#1~\textbf{then}}

\algblockdefx[LOCK]{Trylock}{Untrylock}[1]{\textbf{trylock}~#1~\textbf{then}}

\algblockdefx[VARIABLES]{Variables}{EndVariables}{\textbf{Variables}}

% start with some helper code
% This is the vertical rule that is inserted
\newcommand*{\algrule}[1][\algorithmicindent]{
    \makebox[#1][l]{
        \hspace*{.2em}% <------------- This is where the rule starts from
            \vrule height .75\baselineskip depth .25\baselineskip
    }
}

\newcount\ALG@printindent@tempcnta
\def\ALG@printindent{%
    \ifnum \theALG@nested>0% is there anything to print
        \ifx\ALG@text\ALG@x@notext% is this an end group without any text?
        % do nothing
        \else
        \unskip
        % draw a rule for each indent level
        \ALG@printindent@tempcnta=1
        \loop
        \algrule[\csname ALG@ind@\the\ALG@printindent@tempcnta\endcsname]%
        \advance \ALG@printindent@tempcnta 1
        \ifnum \ALG@printindent@tempcnta<\numexpr\theALG@nested+1\relax
        \repeat
        \fi
        \fi
}
% the following line injects our new indent handling code in place of the default spacing
\patchcmd{\ALG@doentity}{\noindent\hskip\ALG@tlm}{\ALG@printindent}{}{\errmessage{failed to patch    }}
\patchcmd{\ALG@doentity}{\item[]\nointerlineskip}{}{}{} % no spurious vertical space
% end vertical rule patch for algorithmicx

\makeatother
